% GENERATED FILE. Edit canonical lessons, then run npm run generate:books.
% canonical-source-hash: fnv1a64:b12a795ea8ffc68a
% canonical-lessons: BN-C455-dhoirjo, BN-C455-utsaho, BN-C455-agroho, BN-C455-utsuk, BN-C455-tor-sooya

\chapter{Patience and Eagerness}
\label{ch:bn-patience-and-eagerness}
\hlchaptermodality{455}

\begin{chapteropening}
\textbf{By the end of this chapter:} \emph{I can talk about patience and about being eager for something.}

Complete the last lesson of chapter 455: I can talk about patience and about being eager for something.
\end{chapteropening}

\section[\texorpdfstring{\bn{ধৈর্য} (dhoirjô)}{ধৈর্য (dhoirjô)}]{\bn{ধৈর্য} (dhoirjô) — patience}

\label{lesson:BN-C455-dhoirjo}



\begin{quote}
Before the new one: say the Bengali for huddled up, then the Bengali for to hesitate.
\end{quote}

\subsection*{You'll want to know: \bn{ধৈর্য}}
\textbf{\bn{ধৈর্য}} — \emph{dhoirjô} — \textquotedblleft{}patience\textquotedblright{}.

\begin{quote}
\bn{শিশুদের} \bn{পড়াতে} \bn{ধৈর্য} \bn{লাগে।}
\end{quote}

\begin{quote}
\emph{shishuder pôṛāte dhoirjô lāge.}
\end{quote}

\begin{quote}
Teaching children takes patience.
\end{quote}

\bn{ধৈর্য} \bn{ধরা}, to be patient; \bn{ধৈর্য} \bn{হারানো}, to lose patience.

\begin{cousinweb}
\textbf{\bn{ধৈর্য}} (\emph{dhoirjô}) is a tatsama, Sanskrit \emph{dhairya}, \textquotedblleft{}firmness, constancy, calmness, patience\textquotedblright{} (Monier-Williams).
\end{cousinweb}

\begin{grammarlens}[title={What you've built}]
One more word for describing what you lived and felt.
\end{grammarlens}

\subsection*{Your turn}
\begin{itemize}
\raggedright
  \item \emph{Say it:} \emph{dhoirjô}
  \item \emph{Say it:} \emph{dhoirjô}, once more
  \item \emph{Say it:} \emph{shishuder pôṛāte dhoirjô lāge.} — Teaching children takes patience.
  \item \emph{Recall:} say the Bengali for huddled up, then the Bengali for to hesitate, then say \emph{dhoirjô} again
\end{itemize}

\begin{tcolorbox}[breakable,colback=teal!4,colframe=teal!35!black,title={Before you move on}]
What does \bn{ধৈর্য} mean? (\textquotedblleft{}Patience\textquotedblright{}.) Say it once more.
\end{tcolorbox}
\section[\texorpdfstring{\bn{উৎসাহ} (utsāhô)}{উৎসাহ (utsāhô)}]{\bn{উৎসাহ} (utsāhô) — enthusiasm; encouragement}

\label{lesson:BN-C455-utsaho}



\begin{quote}
Before the new one: say the Bengali for to hesitate, then the Bengali for patience.
\end{quote}

\subsection*{You'll want to know: \bn{উৎসাহ}}
\textbf{\bn{উৎসাহ}} — \emph{utsāhô} — \textquotedblleft{}enthusiasm; encouragement\textquotedblright{}.

\begin{quote}
\bn{বাবার} \bn{উৎসাহে} \bn{মীরা} \bn{গান} \bn{শিখল।}
\end{quote}

\begin{quote}
\emph{bābār utsāhe mīrā gān shikhlô.}
\end{quote}

\begin{quote}
With her father's encouragement Mira learned to sing.
\end{quote}

\bn{উৎসাহ} \bn{দেওয়া}, to encourage.

\begin{cousinweb}
\textbf{\bn{উৎসাহ}} (\emph{utsāhô}) is a tatsama, Sanskrit \emph{utsāha}, \textquotedblleft{}power; strength of will, resolution\textquotedblright{} (Monier-Williams).
\end{cousinweb}

\begin{grammarlens}[title={What you've built}]
One more word for describing what you lived and felt.
\end{grammarlens}

\subsection*{Your turn}
\begin{itemize}
\raggedright
  \item \emph{Say it:} \emph{utsāhô}
  \item \emph{Say it:} \emph{utsāhô}, once more
  \item \emph{Say it:} \emph{bābār utsāhe mīrā gān shikhlô.} — With her father's encouragement Mira learned to sing.
  \item \emph{Recall:} say the Bengali for to hesitate, then the Bengali for patience, then say \emph{utsāhô} again
\end{itemize}

\begin{tcolorbox}[breakable,colback=teal!4,colframe=teal!35!black,title={Before you move on}]
What does \bn{উৎসাহ} mean? (\textquotedblleft{}Enthusiasm; encouragement\textquotedblright{}.) Say it once more.
\end{tcolorbox}
\section[\texorpdfstring{\bn{আগ্রহ} (āgrôhô)}{আগ্রহ (āgrôhô)}]{\bn{আগ্রহ} (āgrôhô) — keen interest; eagerness}

\label{lesson:BN-C455-agroho}



\begin{quote}
Before the new one: say the Bengali for patience, then the Bengali for enthusiasm.
\end{quote}

\subsection*{You'll want to know: \bn{আগ্রহ}}
\textbf{\bn{আগ্রহ}} — \emph{āgrôhô} — \textquotedblleft{}keen interest; eagerness\textquotedblright{}.

\begin{quote}
\bn{ইতিহাসে} \bn{অরুণের} \bn{আগ্রহ} \bn{আছে।}
\end{quote}

\begin{quote}
\emph{itihāse ôruner āgrôhô āchhe.}
\end{quote}

\begin{quote}
Arun has a keen interest in history.
\end{quote}

\bn{আগ্রহী}, keen, interested.

\begin{cousinweb}
\textbf{\bn{আগ্রহ}} (\emph{āgrôhô}) is a tatsama, Sanskrit \emph{āgraha}, \textquotedblleft{}insisting on, strong inclination\textquotedblright{} (Monier-Williams); Bengali softened the insistence to keen interest.
\end{cousinweb}

\begin{grammarlens}[title={What you've built}]
One more word for describing what you lived and felt.
\end{grammarlens}

\subsection*{Your turn}
\begin{itemize}
\raggedright
  \item \emph{Say it:} \emph{āgrôhô}
  \item \emph{Say it:} \emph{āgrôhô}, once more
  \item \emph{Say it:} \emph{itihāse ôruner āgrôhô āchhe.} — Arun has a keen interest in history.
  \item \emph{Recall:} say the Bengali for patience, then the Bengali for enthusiasm, then say \emph{āgrôhô} again
\end{itemize}

\begin{tcolorbox}[breakable,colback=teal!4,colframe=teal!35!black,title={Before you move on}]
What does \bn{আগ্রহ} mean? (\textquotedblleft{}Keen interest; eagerness\textquotedblright{}.) Say it once more.
\end{tcolorbox}
\section[\texorpdfstring{\bn{উৎসুক} (utsuk)}{উৎসুক (utsuk)}]{\bn{উৎসুক} (utsuk) — eager; curious to know}

\label{lesson:BN-C455-utsuk}



\begin{quote}
Before the new one: say the Bengali for enthusiasm, then the Bengali for keen interest.
\end{quote}

\subsection*{You'll want to know: \bn{উৎসুক}}
\textbf{\bn{উৎসুক}} — \emph{utsuk} — \textquotedblleft{}eager; curious to know\textquotedblright{}.

\begin{quote}
\bn{ফলাফল} \bn{জানার} \bn{জন্য} \bn{ছাত্ররা} \bn{উৎসুক} \bn{হয়ে} \bn{ছিল।}
\end{quote}

\begin{quote}
\emph{phôlāphôl jānār jônyô chhātrôrā utsuk hôye chhilô.}
\end{quote}

\begin{quote}
The students were eager to know the results.
\end{quote}

\begin{cousinweb}
\textbf{\bn{উৎসুক}} (\emph{utsuk}) is a tatsama, Sanskrit \emph{utsuka}, \textquotedblleft{}restless, uneasy\textquotedblright{} and \textquotedblleft{}anxiously eager\textquotedblright{} (Monier-Williams).
\end{cousinweb}

\begin{grammarlens}[title={What you've built}]
One more word for describing what you lived and felt.
\end{grammarlens}

\subsection*{Your turn}
\begin{itemize}
\raggedright
  \item \emph{Say it:} \emph{utsuk}
  \item \emph{Say it:} \emph{utsuk}, once more
  \item \emph{Say it:} \emph{phôlāphôl jānār jônyô chhātrôrā utsuk hôye chhilô.} — The students were eager to know the results.
  \item \emph{Recall:} say the Bengali for enthusiasm, then the Bengali for keen interest, then say \emph{utsuk} again
\end{itemize}

\begin{tcolorbox}[breakable,colback=teal!4,colframe=teal!35!black,title={Before you move on}]
What does \bn{উৎসুক} mean? (\textquotedblleft{}Eager; curious to know\textquotedblright{}.) Say it once more.
\end{tcolorbox}
\section[\texorpdfstring{\bn{তর} \bn{সওয়া} (tôr sôoyā)}{তর সওয়া (tôr sôoyā)}]{\bn{তর} \bn{সওয়া} (tôr sôoyā) — to be able to wait (mostly with not: can't wait)}

\label{lesson:BN-C455-tor-sooya}



\begin{quote}
Before the new one: say the Bengali for keen interest, then the Bengali for eager.
\end{quote}

\subsection*{You'll want to know: \bn{তর} \bn{সওয়া}}
\textbf{\bn{তর} \bn{সওয়া}} — \emph{tôr sôoyā} — \textquotedblleft{}to be able to wait (mostly with not: can't wait)\textquotedblright{}.

\begin{quote}
\bn{ছুটি} \bn{আসছে}, \bn{মীরার} \bn{আর} \bn{তর} \bn{সইছে} \bn{না।}
\end{quote}

\begin{quote}
\emph{chhuṭi āschhe, mīrār ār tôr sôichhe nā.}
\end{quote}

\begin{quote}
The holidays are coming and Mira can't wait any longer.
\end{quote}

Past (he/she …): \bn{সেই} \bn{ছেলের} \bn{তর} \bn{সইল} \bn{না}, \emph{sei chheler tôr sôilô nā}.

The one who waits is in the genitive: \bn{আমার} \bn{তর} \bn{সইছে} \bn{না}, I can't wait.

\begin{cousinweb}
\textbf{\bn{তর} \bn{সওয়া}} (\emph{tôr sôoyā}) is \emph{tôr}, \textquotedblleft{}waiting\textquotedblright{}, whose origin is not settled, with \emph{sôoyā}, \textquotedblleft{}to bear, to put up with\textquotedblright{} (CDIAL 13304).
\end{cousinweb}

\begin{grammarlens}[title={What you've built}]
One more word for describing what you lived and felt.
\end{grammarlens}

\subsection*{Your turn}
\begin{itemize}
\raggedright
  \item \emph{Say it:} \emph{tôr sôoyā}
  \item \emph{Say it:} \emph{tôr sôoyā}, once more
  \item \emph{Say it:} \emph{chhuṭi āschhe, mīrār ār tôr sôichhe nā.} — The holidays are coming and Mira can't wait any longer.
  \item \emph{Recall:} say the Bengali for keen interest, then the Bengali for eager, then say \emph{tôr sôoyā} again
\end{itemize}

\begin{tcolorbox}[breakable,colback=teal!4,colframe=teal!35!black,title={Before you move on}]
What does \bn{তর} \bn{সওয়া} mean? (\textquotedblleft{}To be able to wait (mostly with not: can't wait)\textquotedblright{}.) Say it once more.
\end{tcolorbox}
